\documentclass[11pt]{article}
\usepackage[margin=1in]{geometry}
\usepackage{amsmath,amssymb,amsthm,mathtools}
\usepackage[colorlinks=true,allcolors=blue]{hyperref}
\setlength{\emergencystretch}{2em}
\newtheorem{theorem}{Theorem}
\newtheorem{proposition}[theorem]{Proposition}
\newtheorem{lemma}[theorem]{Lemma}
\newcommand{\R}{\mathbb R}
\newcommand{\C}{\mathbb C}
\newcommand{\N}{\mathbb N}
\newcommand{\Pol}{\mathbb C[x]}
\newcommand{\diag}{\operatorname{diag}}
\newcommand{\Ran}{\operatorname{Ran}}
\title{Exact native parity transport into the Pollaczek pair\\
and certified finite source--observation bounds}
\author{}
\date{20 September 2026}
\begin{document}
\maketitle

\section{The unchanged arithmetic weight and the literature input}

Retain the original quartet, divisor, and completed numerator
\[
 \rho=\tfrac12+\delta+i\gamma,\quad 0<\delta<\tfrac12,\quad\gamma>2,
 \quad h(s)=\prod_{z\in\{\rho,\bar\rho,1-\rho,1-\bar\rho\}}(s-z)^m,
 \quad m\geq1,\quad g=2\xi,
\]
\[
 w_h(y)=\frac{|(g/h)(1/2+iy)|^2}{2\pi},\quad
 m_k=w_h^{*k},\quad k=4l+1\geq9,\quad
 \mu_h=\int_\R w_h(y)\,dy.
\]
All multiplicities are the original complete multiplicities. No
existence of an off-line quartet is asserted here. Conjugation
symmetry of $g$ and of the root set of $h$ gives $w_h(-y)=w_h(y)$,
and convolution gives the same symmetry for $m_k$.
The source variable, centre and quotient polynomial remain
\[
 S=c+iy,\quad c=k/2,\quad q=[1+k(m-1)](k+1)^2,
\]
\[
 \chi(S)=\prod_{a,b=0}^k
 [S-c-(2a-k)\delta-i(2b-k)\gamma]^{1+k(m-1)}.
\]
All finite observation formulas below use the original degree
$N\geq q-1$ and actual map $\mathcal A_S=\Lambda\mathcal J_{S,N}$,
with its full physical Taylor unit. Its codomain is its original
observation image with the retained Hermitian structure.

The original proved CAI16--18 envelope is used at its actual scope
\cite{CAI}%
% BEGIN VERIFIED PUBLIC PROOF LINK
~\href{https://github.com/KokunoYumeto/zeta-function-research-reader/blob/487253b5b01851da5b9b98af4e146fecf1fab902/workbenches/splitzero-tandem/continuations/20260919-full-window-source-products/providers/CAI_COMPLETE_PREREQUISITE.tex\#L284}{[proof CAI16]}%
% END VERIFIED PUBLIC PROOF LINK
{}; the arbitrary-fixed-multiplicity form is HG21 \cite{HG}%
% BEGIN VERIFIED PUBLIC PROOF LINK
~\href{https://github.com/KokunoYumeto/zeta-function-research-reader/blob/5992ad50c0f2a06921f1fcaded7b4e38097c0739/workbenches/splitzero-tandem/continuations/20260920-original-kernel-mixed-determinants/providers/HEAT_GROWTH_AND_ORIGINAL_CURRENTS.tex\#L374}{[proof HG21]}%
% END VERIFIED PUBLIC PROOF LINK
{}:
\begin{equation}\label{eq:CAI}
 a_k(1+y^2)^{-M}\rho_\sigma(y)\leq m_k(y)\leq A_k\rho_\sigma(y),
 \quad M=21+4m,
\end{equation}
\[
 \rho_\sigma(y)=\frac{|\Gamma(1/4+iy/2)|^2}{2\pi},\quad
 \int\rho_\sigma=\sqrt{2\pi},\quad\alpha=\pi/2,
\]
\[
 c_\Gamma e^{-\alpha|y|}(1+|y|)^{-1/2}
 \leq\rho_\sigma(y)\leq
 C_\Gamma e^{-\alpha|y|}(1+|y|)^{-1/2},
\]
\[
 c_\Gamma=\sqrt2e^{-7/6},\quad C_\Gamma=\sqrt{2\sqrt5}e^{2/3},
 \quad p=8m-9/2,\quad B_h=42+8m,
\]
\[
 a_k=\frac{c_h\vartheta_h^{k-3}e^{-\alpha(k-3)}(k-2)^{-B_h}}
                   {C_\Gamma2^{B_h/2}},\quad
 A_k=\frac{C_h^{\rm tilt}}{c_\Gamma}k^{p+1}M_\alpha^{k-1},
\]
\[
 \vartheta_h=\int_{-1}^1w_h(y)\,dy,\qquad
 M_\alpha=\int_\R e^{\alpha y}w_h(y)\,dy.
\]
The constants $c_h,C_h^{\rm tilt}$ are precisely the established
lower-envelope and compact/tail constants in the complete CAI
provider, not new assumptions. In particular $m_k>0$ and all
moments exist. For complex polynomials $P$ of degree at most $d$,
that provider proves
\begin{equation}\label{eq:CAI-finite}
 \frac{a_k}{\Delta_d}\|P\|_\sigma^2
       \leq\|P\|_{m_k}^2\leq A_k\|P\|_\sigma^2,
 \qquad \Delta_d=[1+4(d+M)^2]^M.
\end{equation}
The Gamma mass has not been divided out.

The original author source used here is A.~I.~Aptekarev,
G.~L\'opez Lagomasino and A.~Mart\'{\i}nez--Finkelshtein,
\emph{Strong asymptotics for the Pollaczek multiple orthogonal
polynomials ensembles},
\href{https://arxiv.org/abs/1410.1261v1}{arXiv:1410.1261v1}
\cite{ALM}.
Its Introduction gives the weights
\begin{equation}\label{eq:literature}
 w_1(x)=\frac1{\sinh(\alpha\sqrt x)},\qquad
 w_2(x)=\frac1{\sqrt x\cosh(\alpha\sqrt x)},\qquad x>0,
\end{equation}
and their positive Stieltjes ratio
\begin{equation}\label{eq:S}
 s(x)=\frac{w_2(x)}{w_1(x)}
     =\frac{\tanh(\alpha\sqrt x)}{\sqrt x}
     =\frac4\pi\sum_{j=0}^{\infty}\frac1{x+(2j+1)^2}.
\end{equation}
The lowercase $s(x)$ in this note denotes this multiplier, not
the original complex source variable $S$. No multiple-orthogonal
asymptotic is assumed for the native weight.

\section{Exact parity-square maps, source images, and masses}

Define the two original transported measures on $(0,\infty)$ by
\begin{equation}\label{eq:nu}
 d\nu_0(x)=\frac{m_k(\sqrt x)}{\sqrt x}\,dx,
 \qquad d\nu_1(x)=\sqrt x\,m_k(\sqrt x)\,dx=x\,d\nu_0(x).
\end{equation}
They have strictly positive densities. Their literal density
ratios against the first literature weight are
\begin{equation}\label{eq:ratios}
 r_0(x)=\frac{m_k(\sqrt x)\sinh(\alpha\sqrt x)}{\sqrt x},
 \qquad r_1(x)=x r_0(x),\qquad d\nu_j=r_jw_1\,dx.
\end{equation}
For every polynomial $f(y)$ there is a unique pair $E,O$ with
$f(y)=E(y^2)+yO(y^2)$. Evenness of $m_k$ makes the cross term in
the squared norm odd. The substitution $x=y^2$ on each half-axis
then proves the exact sesquilinear and norm identities
\begin{equation}\label{eq:parity-norm}
 \|f\|_{m_k}^2=\int_0^\infty|E(x)|^2\,d\nu_0(x)
                    +\int_0^\infty|O(x)|^2\,d\nu_1(x).
\end{equation}
There is no missing factor two: the two half-axes cancel the
$1/2$ in $dy=dx/(2\sqrt x)$.

The same construction is a unitary isomorphism
\[
 L^2(\R,m_kdy)\longrightarrow L^2(\nu_0)\oplus L^2(\nu_1),
\]
with
\[
 E(x)=\frac{f(\sqrt x)+f(-\sqrt x)}2,\qquad
 O(x)=\frac{f(\sqrt x)-f(-\sqrt x)}{2\sqrt x}.
\]
These formulas are defined almost everywhere, and their inverse
is $f(y)=E(y^2)+yO(y^2)$ for $y\ne0$. The identity of norms proves
both that the maps are well defined and that they are inverse
isometries.

The exact isometry into the literature's two Hilbert spaces is
\begin{equation}\label{eq:literature-isometry}
 (E,O)\longmapsto
       \left(\sqrt{r_0}\,E,\sqrt{r_1/s}\,O\right)
 \quad\text{in }L^2(w_1dx)\oplus L^2(w_2dx).
\end{equation}
Since all density ratios are positive, division by their square
roots gives its inverse on these full Hilbert spaces. On the
degree-$N$ source its image is exactly
\[
 \sqrt{r_0}\,\Pol_{\lfloor N/2\rfloor}
       \ \oplus\ \sqrt{r_1/s}\,\Pol_{\lfloor(N-1)/2\rfloor},
\]
not an unweighted polynomial pair. Here $\Pol_n$ means degree at
most $n$, and $\Pol_{-1}=\{0\}$. The common-first-weight version
is $(E,O)\mapsto(\sqrt{r_0}E,\sqrt{r_1}O)$ in two copies of
$L^2(w_1dx)$.

Write $\eta_j=\int y^jw_h(y)\,dy$ and
$\zeta_j=\int y^jm_k(y)\,dy$. Tonelli and the convolution
multinomial formula give the original parity masses
\begin{equation}\label{eq:masses}
 \nu_0((0,\infty))=\zeta_0=\mu_h^k,
 \qquad \nu_1((0,\infty))=\zeta_2=k\eta_2\mu_h^{k-1}.
\end{equation}
The cross terms in the second formula vanish because $w_h$ is
even. The literature masses are different fixed numbers:
\begin{equation}\label{eq:raw-moments}
 \tau_j:=\int_0^\infty x^j w_1(x)\,dx
 =\frac{4(2j+1)!}{\alpha^{2j+2}}
        (1-2^{-2j-2})\zeta(2j+2),\qquad \tau_0=2,
\end{equation}
\[
 \int_0^\infty w_2(x)\,dx=2.
\]
For the first formula use $x=y^2$ and the positive expansion
$1/\sinh(\alpha y)=2\sum_{j\geq0}e^{-(2j+1)\alpha y}$,
then integrate each power of $y$. Its constant case uses
$\sum_{j\geq0}(2j+1)^{-2}=\pi^2/8$, also obtained below from
\eqref{eq:S} at zero. For the second mass, the same substitution
gives $2\int_0^\infty\operatorname{sech}(\alpha y)dy=\pi/\alpha=2$;
an antiderivative of $\operatorname{sech} u$ is
$\arctan(\sinh u)$. None of these masses is set equal to one.

For the exact original source $F(S)=\sum_{j=0}^Na_jS^j$, define
\[
 T_{rj}=\binom jr c^{j-r}i^r\quad(r\leq j),\qquad T_{rj}=0\quad(r>j).
\]
Let $\Pi_N$ reorder the coefficients in powers of $y$ into even
powers first and odd powers second, without rescaling any entry.
The parity coefficient map is the invertible matrix
\begin{equation}\label{eq:source-map}
 L_N=\Pi_NT,\qquad (e,o)=L_Na,
\end{equation}
where explicitly
\[
 e_r=(-1)^r\sum_{j=2r}^N\binom j{2r}c^{j-2r}a_j,
 \quad
 o_r=i(-1)^r\sum_{j=2r+1}^N\binom j{2r+1}c^{j-2r-1}a_j.
\]
Thus the full original degree-$N$ source maps onto precisely the
displayed pair of polynomial spaces. An original constrained
source subspace $W$ maps to $L_NW$, not to a substituted full pair.

Put $n_0=\lfloor N/2\rfloor$, $n_1=\lfloor(N-1)/2\rfloor$ and
\begin{equation}\label{eq:parity-grams}
 H_{j,n}=\left[\int_0^\infty x^{a+b}\,d\nu_j(x)\right]_{a,b=0}^n,
 \qquad H=\diag(H_{0,n_0},H_{1,n_1}).
\end{equation}
The matrices are positive definite because nonzero polynomials
cannot vanish on a set of full measure for a positive density.
They retain the literal moments
$(H_{j,n})_{ab}=\zeta_{2(a+b+j)}$. Integration gives
\begin{equation}\label{eq:original-gram}
 H_{S,N}^{\rm ar}=L_N^*HL_N.
\end{equation}
The actual observation matrix in these coordinates is
\begin{equation}\label{eq:observation}
A=\mathcal A_SL_N^{-1}.
\end{equation}
The exact original observation metric and minimizing lift are
\begin{equation}\label{eq:Q}
 Q(H)=(AH^{-1}A^*)^{-1},\qquad
 a_{\min}=L_N^{-1}H^{-1}A^*Q(H)u.
\end{equation}
Indeed $A$ is onto its stated observation image, so $AH^{-1}A^*$
is strictly positive. The vector $z_0=H^{-1}A^*Q(H)u$ satisfies
$Az_0=u$, and for $Az=0$ the cross term is
$z^*Hz_0=(Az)^*Q(H)u=0$. Every other lift has energy
$u^*Q(H)u+z^*Hz$, proving the formula and uniqueness. The
congruence \eqref{eq:original-gram} returns the result to the
unchanged source coefficients.

There is also an exact transport of the original quotient, not
only of its source norm. Put
\[
 u_a=(2a-k)\delta,\qquad v_b=(2b-k)\gamma,\qquad
 e_k=1+k(m-1),
\]
and define the monic polynomial
\begin{equation}\label{eq:Xi}
 \Xi(x)=\prod_{a=0}^k\prod_{b=(k+1)/2}^{k}
                 [x-(v_b-iu_a)^2]^{e_k}.
\end{equation}
Its degree is $q/2$. The opposite grid point of $(a,b)$ is
$(k-a,k-b)$, and exactly one point of each opposite pair has
$v_b>0$. That pair's factors after $S=c+iy$ are
\[
 (iy-u_a-iv_b)(iy+u_a+iv_b)
          =-[y^2-(v_b-iu_a)^2].
\]
Since $k$ is odd, $q=e_k(k+1)^2$ is divisible by four. The number
$q/2$ of paired factors with multiplicity is therefore even,
so the full sign is $(-1)^{q/2}=1$. Consequently
\begin{equation}\label{eq:chi-square}
 \chi(c+iy)=\Xi(y^2)
\end{equation}
with exactly the original leading coefficient. Conjugating the
roots in \eqref{eq:Xi} is the permutation $a\mapsto k-a$, so
$\Xi$ also has real coefficients.

Let $\mathcal E_x=\C[x]/(\Xi)$. Substitution and parity give the
explicit algebra isomorphism followed by vector-space coordinates
\begin{equation}\label{eq:quotient-parity}
 \C[S]/(\chi)\longrightarrow\C[y]/(\Xi(y^2))
                  \longrightarrow\mathcal E_x\oplus\mathcal E_x,
 \qquad [F]\longmapsto([E],[O]).
\end{equation}
The second arrow has multiplication
$(E,O)(E',O')=(EE'+xOO',EO'+OE')$ modulo $\Xi$ in each component.
To verify its kernel, divide $E$ and $O$ separately by $\Xi$.
The remaining even and odd polynomials have degrees less than
$q$ in $y$, so their sum is divisible by the monic degree-$q$
polynomial $\Xi(y^2)$ only when both remainders are zero. This
also proves surjectivity and identifies the inverse through the
unique remainder $E(y^2)+yO(y^2)$ and $y=(S-c)/i$.
Multiplication by the original variable $S$ is precisely
\begin{equation}\label{eq:S-block}
 \begin{pmatrix}cI&iM_x\\ iI&cI\end{pmatrix}
 \quad\text{on }\mathcal E_x\oplus\mathcal E_x.
\end{equation}
An original full Taylor unit $\mathfrak u\in\C[S]/(\chi)$,
transported as $(U_0,U_1)$, acts by the unchanged block matrix
\begin{equation}\label{eq:unit-block}
 \begin{pmatrix}M_{U_0}&M_{xU_1}\\M_{U_1}&M_{U_0}\end{pmatrix}.
\end{equation}
Its inverse is the transported multiplication by
$\mathfrak u^{-1}$; the unit has not been removed. Thus the
original class and observation maps can equivalently be
transported through \eqref{eq:quotient-parity} and
\eqref{eq:unit-block}. Formula \eqref{eq:observation} is their
same matrix composition before choosing these quotient
coordinates. At $N\geq q-1$, both parity source degree bounds
are at least $q/2-1$, so both residue-class components are
attained by the original degree-$N$ source.

\section{Both native orthogonal families and all Christoffel constants}

Let $P_d(y)$ be the real monic orthogonal polynomial of degree $d$
for $m_k(y)dy$ and let $\omega_d=\int P_d^2m_k>0$. Existence and
uniqueness follow by solving the positive definite moment system.
Evenness and uniqueness give $P_d(-y)=(-1)^dP_d(y)$, so define
\[
 P_{2n}(y)=E_n(y^2),\quad P_{2n+1}(y)=yO_n(y^2),
 \quad \kappa_n=\omega_{2n},\quad\lambda_n=\omega_{2n+1}.
\]
Equation \eqref{eq:parity-norm} shows that $E_n$ and $O_n$ are
respectively the real monic orthogonal families for $\nu_0$ and
$\nu_1$, with norms $\kappa_n$ and $\lambda_n$.

The monic orthogonal polynomial in the original complex source
coordinate is not obtained by dropping the phase. It is exactly
\begin{equation}\label{eq:S-polynomial-phase}
 \mathsf P_d(S)=i^dP_d((S-c)/i),\qquad
 \mathsf P_d(c+iy)=i^dP_d(y),\qquad
 \|\mathsf P_d\|_{m_k}^2=\omega_d.
\end{equation}
Its leading coefficient is $i^di^{-d}=1$; the lower-degree
orthogonality follows by the same invertible substitution used
in \eqref{eq:source-map}. Thus uniqueness proves this exact
dictionary. Its parity pair is $((-1)^nE_n,0)$ at degree $2n$
and $(0,i(-1)^nO_n)$ at degree $2n+1$. These phases remain in
the original boundary classes and observation entries.

Orthogonality gives the full original recurrence
\[
 yP_d=P_{d+1}+\beta_dP_{d-1},\qquad
 \beta_d=\omega_d/\omega_{d-1}>0\quad(d\geq1),\quad \beta_0=0.
\]
To prove it, expand $yP_d-P_{d+1}$ in the lower orthogonal basis.
Its coefficients below $d-1$ vanish by transferring $y$ to the
lower polynomial. The coefficient of $P_d$ vanishes by parity,
and the coefficient of $P_{d-1}$ is
$\langle P_d,yP_{d-1}\rangle/\omega_{d-1}=\omega_d/\omega_{d-1}$.
In the original $S$ coordinate the same identity reads
$(S-c)\mathsf P_d=\mathsf P_{d+1}-\beta_d\mathsf P_{d-1}$,
since $i^{d+1}/i^{d-1}=-1$; its sign is therefore retained too.

Set
\begin{equation}\label{eq:ad}
 a_n=\frac{\lambda_n}{\kappa_n}=\beta_{2n+1}>0,
 \qquad d_n=\frac{\kappa_n}{\lambda_{n-1}}=\beta_{2n}>0\ (n\geq1),
 \qquad d_0=0.
\end{equation}
Separating the even and odd instances of that recurrence gives
the exact polynomial identities
\begin{equation}\label{eq:Christoffel}
 xO_n=E_{n+1}+a_nE_n,\qquad
 E_n=O_n+d_nO_{n-1},\quad O_{-1}=0.
\end{equation}
Evaluating the first at zero proves, starting from $E_0=1$,
\begin{equation}\label{eq:E0}
 E_n(0)=(-1)^n\prod_{j=0}^{n-1}a_j\ne0,
 \qquad \frac{E_{n+1}(0)}{E_n(0)}=-a_n.
\end{equation}
Therefore the exact Christoffel formula is
\begin{equation}\label{eq:Christoffel-explicit}
 O_n(x)=\frac{E_{n+1}(x)-[E_{n+1}(0)/E_n(0)]E_n(x)}x
        =\frac{E_{n+1}(x)+a_nE_n(x)}x,
 \quad \lambda_n=a_n\kappa_n.
\end{equation}
The numerator has a zero at zero, so this is a polynomial identity,
including the endpoint. The complete odd zero value is
\begin{equation}\label{eq:O0}
 O_n(0)=(-1)^n\sum_{j=0}^n
     \left(\prod_{r=j+1}^nd_r\right)
     \left(\prod_{r=0}^{j-1}a_r\right)
 =E_{n+1}'(0)+a_nE_n'(0).
\end{equation}
The sum follows by iterating $O_n(0)=E_n(0)-d_nO_{n-1}(0)$;
the derivative formula follows from \eqref{eq:Christoffel-explicit}.
Every empty product equals one, so this covers $n=0$.

Substitution of the two identities in \eqref{eq:Christoffel} into
each other yields both complete monic recurrences
\begin{align}
 xE_n&=E_{n+1}+(a_n+d_n)E_n+a_{n-1}d_nE_{n-1},\label{eq:E-rec}\\
 xO_n&=O_{n+1}+(a_n+d_{n+1})O_n+a_nd_nO_{n-1}.\label{eq:O-rec}
\end{align}
The last term of the first is zero at $n=0$ and needs no $a_{-1}$.
The norm ratios are exactly
$\kappa_n/\kappa_{n-1}=a_{n-1}d_n$ and
$\lambda_n/\lambda_{n-1}=a_nd_n$.
For example the orthonormal version of the first cross relation is
\[
 x\frac{O_n}{\sqrt{\lambda_n}}
  =\sqrt{d_{n+1}}\frac{E_{n+1}}{\sqrt{\kappa_{n+1}}}
       +\sqrt{a_n}\frac{E_n}{\sqrt{\kappa_n}},
\]
and the second is
\[
 \frac{E_n}{\sqrt{\kappa_n}}
 =\sqrt{a_n}\frac{O_n}{\sqrt{\lambda_n}}
       +\sqrt{d_n}\frac{O_{n-1}}{\sqrt{\lambda_{n-1}}}.
\]
These equalities are between functions in the displayed different
weighted spaces; they explicitly give the matrix of the identity
map or $x$-multiplication on every finite polynomial domain.

\section{The exact multiple-orthogonality receiver and its residuals}

The raw paper has one polynomial $P$ subject to moments against
both $w_1$ and $w_2$. The native common-multiplier Nikishin pair is
\[
 (d\nu_0,\ s\,d\nu_0)=r_0(x)(w_1(x),w_2(x))\,dx;
\]
the native odd measure is instead $d\nu_1=x\,d\nu_0$.
The literal maps between their functionals, for every polynomial
$P$ and nonnegative integer $j$, are
\begin{align}
 \int x^jP\,d\nu_0&=\int x^j(r_0P)w_1\,dx,\label{eq:functional0}\\
 \int x^jP\,s\,d\nu_0&=\int x^j(r_0P)w_2\,dx,\label{eq:functionalS}\\
 \int x^jP\,d\nu_1&=\int x^j\left(\frac{r_1}{s}P\right)w_2\,dx
                    =\int x^{j+1}P\,d\nu_0.\label{eq:functional1}
\end{align}
Thus each finite raw moment operator receives the explicit
multiplied source space in these formulas, not the raw polynomial
space unless that multiplier is retained. These maps are the
functional counterparts of the square-root isometry
\eqref{eq:literature-isometry}.

Here is a nonzero residual computation within the native common
multiplier pair. The $n$ zeros of $E_n$ are simple and lie in
$(0,\infty)$. To prove this, multiply $E_n$ by the product of its
distinct sign-changing zeros in that interval. If fewer than $n$
such zeros existed, this product would have degree less than $n$
and its product with $E_n$ would have constant nonzero sign except
at finitely many points of the positive interval. Positivity of
$\nu_0$ would contradict orthogonality. Thus there are $n$ such
zeros, which exhaust the degree and establish the assertion.
In particular $E_n(-a)$ has sign $(-1)^n$ for every $a>0$.

For $a>0$, divide $E_n(x)-E_n(-a)$ by $x+a$. The quotient has
degree $n-1$, so orthogonality gives the exact positive-resolvent
identity
\begin{equation}\label{eq:residual-atom}
 \int\frac{E_n(x)}{x+a}\,d\nu_0(x)
  =\frac1{E_n(-a)}\int\frac{E_n(x)^2}{x+a}\,d\nu_0(x).
\end{equation}
For $0\leq j<n$, division of $x^j$ by $x+a$ similarly gives
\begin{equation}\label{eq:residual}
 \int x^jE_n(x)s(x)\,d\nu_0(x)
 =\frac4\pi\sum_{r=0}^{\infty}
 \frac{[-(2r+1)^2]^j}{E_n(-(2r+1)^2)}
 \int\frac{E_n(x)^2}{x+(2r+1)^2}\,d\nu_0(x).
\end{equation}
To justify the sum before division, dominate by
$\int|x^jE_n(x)|s(x)\,d\nu_0<\infty$; the atomwise integrals are
absolutely summable since $s$ is bounded. The divisions preserve
each integral exactly. Every term on the resulting right side
has the same strict sign $(-1)^{n+j}$, so that is the sign of the
nonzero residual. In particular $E_n$ does not satisfy even the
zeroth second-weight condition of the native common-multiplier
Nikishin system. This is an evaluated residual, not just a
declaration that two orthogonalities differ. For $n=0$ the
zeroth residual is directly positive.

There is also a completely finite lower bound. Cauchy--Schwarz
with the factors $(x+1)^{1/2}E_n$ and $(x+1)^{-1/2}E_n$ yields
\[
 \int\frac{E_n^2}{x+1}\,d\nu_0
 \geq\frac{\kappa_n^2}{\int(x+1)E_n^2\,d\nu_0}
 =\frac{\kappa_n}{1+a_n+d_n},
\]
where \eqref{eq:E-rec} computes the final moment. The first
positive atom in \eqref{eq:residual} therefore gives
\begin{equation}\label{eq:residual-lower}
 \left|\int E_n s\,d\nu_0\right|
 \geq\frac{4\kappa_n}{\pi(1+a_n+d_n)|E_n(-1)|}>0.
\end{equation}
All its constants belong to the original native parity family.

\section{Two positive resolvent series, with all constants proved}

We give a direct proof of \eqref{eq:S} and the reciprocal identity
that will supply finite native bounds. On the interval
$[0,\alpha]$, $\alpha=\pi/2$, the operator $-d^2/dt^2+x$ with
Dirichlet condition at zero and Neumann condition at $\alpha$
has eigenvalues $x+(2j+1)^2$ and normalized eigenfunctions
$\sqrt{2/\alpha}\sin((2j+1)t)$. Their completeness is the ordinary
sine Fourier expansion obtained by odd reflection at zero and
even reflection at $\alpha$. The Green kernel is
\[
 G_x(u,v)=\frac{\sinh(\sqrt x\min(u,v))
                  \cosh(\sqrt x(\alpha-\max(u,v)))}
                    {\sqrt x\cosh(\alpha\sqrt x)}.
\]
Its formula follows by solving the homogeneous differential
equation on either side of $u=v$, imposing the two boundary
conditions, continuity, and the unit derivative jump. Its
eigenfunction expansion is uniformly absolutely convergent on
the closed square since its terms are bounded by
$(2/\alpha)/(x+(2j+1)^2)$. The spectral resolvent and the displayed
Green kernel agree in $L^2$ and hence by continuity everywhere.
At $u=v=\alpha$ this gives \eqref{eq:S}; at $x=0$, continuity
gives $s(0)=\alpha$ and the odd reciprocal-square sum used above.

With Neumann conditions at both endpoints, the constant mode
contributes $1/(\alpha x)$ and the positive modes have
eigenvalues $x+4j^2$, $j\geq1$. The normalized endpoint values
have squares $2/\alpha$. The analogous Green kernel has endpoint
value $\coth(\alpha\sqrt x)/\sqrt x$. Multiplication by $x$ gives
the exact reciprocal series
\begin{equation}\label{eq:reciprocal}
 \boxed{\frac1{s(x)}=\sqrt x\coth(\alpha\sqrt x)
  =\frac2\pi+\frac4\pi\sum_{j=1}^{\infty}\frac{x}{x+4j^2}.}
\end{equation}
The same uniformly convergent kernel argument justifies the
expansion for $x>0$. Both sides extend at zero to $2/\pi$.
Every summand is nonnegative. These identities are used on the
original axis $x>0$; no degree-dependent rescaling is made.

\section{Finite Stieltjes compression and native Gram approximation}

The following formulas apply to each original parity block with
$\nu=\nu_j$ and degree $n=n_j$. Put $v_n(x)=(1,x,\ldots,x^n)^T$,
\[
 H=\int v_nv_n^*\,d\nu,\quad
 X_H=\int x v_nv_n^*\,d\nu,\quad
 K=\int v_nv_n^*s\,d\nu,\quad
 K_x=\int x v_nv_n^*s\,d\nu.
\]
All four are finite positive definite matrices, and for the
zeroth parity block $X_H=H_{1,n}$ exactly. For $a>0$ define
$R(a)=\int v_nv_n^*/(x+a)\,d\nu$. The block Gram matrix of
$(x+a)^{1/2}v_n$ and $(x+a)^{-1/2}v_n$ is positive. Its Schur
complement gives
\begin{equation}\label{eq:resolvent-compression}
 R(a)\succeq H(X_H+aH)^{-1}H.
\end{equation}
The positive series \eqref{eq:S} therefore proves
\begin{equation}\label{eq:finite-S}
 H^{1/2}s(H^{-1/2}X_HH^{-1/2})H^{1/2}
       \preceq K\preceq\alpha H.
\end{equation}
The matrix function in the lower bound is defined by the positive
spectral decomposition of the finite matrix. The summation is
justified because each inverse is at most $a^{-1}I$ and the
reciprocal-square series converges. This is a finite-moment
bound using the exact Christoffel moment matrix, not a replacement
of the full multiplication operator by its compression.

More directly useful for the unchanged original Gram is
\eqref{eq:reciprocal}. Define the actual second-weight tilt
$d\widetilde\nu=s\,d\nu=r_jw_2\,dx$ and, for an integer $J\geq1$,
\begin{equation}\label{eq:H-trunc}
 H^{[J]}=\frac2\pi K+\frac4\pi\sum_{j=1}^J
       \int\frac{x}{x+4j^2}v_nv_n^*\,d\widetilde\nu.
\end{equation}
Every summand is positive and $H^{[J]}$ is strictly positive
already from its first term. Integrating \eqref{eq:reciprocal}
and using monotone convergence in every quadratic form gives
$H^{[J]}\uparrow H$. Since
$x/(x+4j^2)\leq x/(4j^2)$ and
$\sum_{j>J}j^{-2}\leq\int_J^\infty t^{-2}dt=1/J$, the full
matrix error bound is
\begin{equation}\label{eq:H-error}
 \boxed{0\preceq H-H^{[J]}\preceq\frac1{\pi J}K_x
                                      \preceq\frac1{2J}X_H.}
\end{equation}
This is a certified approximation of the original native Gram,
not just a representation of an auxiliary tilted metric.

Only moments and negative-axis resolvent integrals are required
for a finite truncation. If
$\widetilde\mu_d=\int x^d\,d\widetilde\nu$ and
$\widetilde F(a)=\int(x+a)^{-1}\,d\widetilde\nu$, polynomial
division gives, for $d\geq1$,
\begin{equation}\label{eq:resolvent-moments}
 \int\frac{x^d}{x+a}\,d\widetilde\nu
 =\sum_{r=0}^{d-1}(-a)^r\widetilde\mu_{d-1-r}
                  +(-a)^d\widetilde F(a).
\end{equation}
For $d=0$ the integral is $\widetilde F(a)$. This proves the
finite evaluation map, including every polynomial-division sign.
The positive integral form in \eqref{eq:H-trunc} remains available
when cancellation in \eqref{eq:resolvent-moments} would be
numerically inconvenient.

Apply these constructions to both blocks of
\eqref{eq:parity-grams}. Denote their direct sums by $H^{[J]}$,
$K_x$ and $X_H$. The original source error is exactly the
congruence
\[
 0\preceq H_{S,N}^{\rm ar}-L_N^*H^{[J]}L_N
       \preceq\frac1{\pi J}L_N^*K_xL_N.
\]
For the unchanged onto observation $A$ in \eqref{eq:observation},
the minimum-energy proof of \eqref{eq:Q} shows that the map
$G\mapsto Q(G)=(AG^{-1}A^*)^{-1}$ is order-preserving on positive
source matrices: taking the minimum over the same set $Az=u$
preserves each quadratic-form inequality. Therefore
\begin{equation}\label{eq:Q-bracket}
 \boxed{Q(H^{[J]})\preceq Q(H)\preceq
 Q\left(H^{[J]}+\frac1{\pi J}K_x\right).}
\end{equation}
Both endpoints converge to the original $Q(H)$: the finite
matrices converge to $H$, their positive inverses converge by
$B^{-1}-C^{-1}=B^{-1}(C-B)C^{-1}$, and the same holds after
the onto observation congruence. This proves a finite certified
source and observation calculation with the original classes and
Taylor unit untouched.

There is an entirely explicit degree-only bound for the error.
The original Gamma orthonormal recurrence has off-diagonal
$b_j=\sqrt{(j+1)(j+1/2)}\leq j+1$, so its two weighted shifts
give $\|yP\|_\sigma\leq2(N+1)\|P\|_\sigma$ on degree at most $N$.
Combining this with \eqref{eq:CAI-finite} yields
\begin{equation}\label{eq:L-bound}
 \|yP\|_{m_k}^2\leq\mathcal L_N\|P\|_{m_k}^2,
 \qquad \mathcal L_N=
       \frac{4A_k\Delta_N(N+1)^2}{a_k}.
\end{equation}
Under the parity map this is exactly $X_H\preceq\mathcal L_NH$.
Consequently for $J>\mathcal L_N/2$, with
$\varepsilon_J=\mathcal L_N/(2J)<1$, one has
\begin{equation}\label{eq:relative}
 (1-\varepsilon_J)H\preceq H^{[J]}\preceq H,
 \qquad
 Q(H^{[J]})\preceq Q(H)
       \preceq(1-\varepsilon_J)^{-1}Q(H^{[J]}).
\end{equation}
This bound is generally coarse but is finite, positive and
explicit in the original CAI constants and degree. It does not
assume the desired arithmetic estimate.

\section{A second fully explicit bound using raw Pollaczek moments}

The literal ratio $r_0$ also admits a useful finite comparison
with the raw first literature weight. Put $L=M+1$ and
\begin{equation}\label{eq:ratio-constants}
 c_0=\frac{a_kc_\Gamma\alpha}{2^{1/4}(1+2\alpha)},
 \qquad C_0=A_kC_\Gamma\alpha.
\end{equation}
Then for every $x>0$,
\begin{equation}\label{eq:ratio-envelope}
 c_0(1+x)^{-L}\leq r_0(x)\leq C_0.
\end{equation}
To prove this, write $y=\sqrt x$ and retain the exact factor
\[
 e^{-\alpha y}\frac{\sinh(\alpha y)}y
    =\frac{1-e^{-2\alpha y}}{2y}.
\]
It is at most $\alpha$ and at least
$\alpha/(1+2\alpha y)$, since $1-e^{-u}\geq u/(1+u)$ for $u\geq0$.
The latter inequality follows from $e^u\geq1+u$.
Substituting the two CAI envelopes gives the upper bound and the
lower expression
\[
 \frac{a_kc_\Gamma\alpha(1+y^2)^{-M}}
                   {(1+2\alpha y)(1+y)^{1/2}}.
\]
Use $1+2\alpha y\leq(1+2\alpha)\sqrt{1+y^2}$ and
$(1+y)^{1/2}\leq2^{1/4}(1+y^2)^{1/4}$, followed by
$(1+y^2)^{-M-3/4}\geq(1+y^2)^{-M-1}$, to prove
\eqref{eq:ratio-envelope} with the stated constants.

For parity $j=0,1$ use the raw measure
$d\tau^{(j)}=x^jw_1(x)dx$ and the degree $n=n_j$. Define the
explicit moment matrices
\[
 U_{j,n}=[\tau_{a+b+j}]_{a,b=0}^n,\qquad
 V_{j,n}=\left[\sum_{r=0}^L\binom Lr\tau_{a+b+j+r}\right]_{a,b=0}^n.
\]
All entries are the explicit values \eqref{eq:raw-moments}.
The measures satisfy $d\nu_j=r_0\,d\tau^{(j)}$, so the same
block-Gram Schur-complement argument used in
\eqref{eq:resolvent-compression}, now with weight $(1+x)^L$,
proves the actual native bound
\begin{equation}\label{eq:raw-matrix-bound}
 \boxed{c_0 U_{j,n}V_{j,n}^{-1}U_{j,n}
             \preceq H_{j,n}\preceq C_0U_{j,n}.}
\end{equation}
Explicitly, the block matrix with blocks $V_{j,n},U_{j,n}$ and
$\int v_nv_n^*(1+x)^{-L}d\tau^{(j)}$ is a Gram matrix, giving
the lower rational-weight bound; \eqref{eq:ratio-envelope}
then gives the displayed native inequality.
Taking block direct sums, followed by the unchanged congruence
$L_N$ or the unchanged observation $A$, transports these bounds
to the original source or to $Q(H)$, respectively. In particular,
if $G_-$ and $G_+$ are these two block bounds, then
\[
 Q(G_-)\preceq Q(H)\preceq Q(G_+).
\]
This gives a fully explicit finite comparison from raw literature
moments to the actual observation metric, with every density
ratio paid for. No raw multiple-orthogonal polynomial is being
used as a native orthogonal polynomial.

The same ratio envelope proves a typed map in the reverse
direction for the paper's actual multiple orthogonal polynomial
$P_{(n_1,n_2)}$. The function
\[
 f(x)=P_{(n_1,n_2)}(x)/r_0(x)
\]
belongs to $L^2(\nu_0)$, since its squared norm is
$\int |P|^2r_0^{-1}w_1dx$ and the inverse ratio is bounded by
$c_0^{-1}(1+x)^L$. Equations
\eqref{eq:functional0}--\eqref{eq:functionalS} give exactly
\[
 \int x^jf\,d\nu_0=0\ (0\leq j<n_1),\qquad
 \int x^jf\,s\,d\nu_0=0\ (0\leq j<n_2).
\]
This is an explicit receiver of the raw multiple-orthogonality
relations in the original even Hilbert sector. Its source image
is $r_0^{-1}\Pol$, not the original polynomial source. For the
paper's diagonal rescaling
$Q_n(x)=(4n^2)^{-2n}P_{(n,n)}(4n^2x)$, the corresponding original
even-axis function is exactly
\[
 y\longmapsto\frac{(4n^2)^{2n}Q_n(y^2/(4n^2))}{r_0(y^2)}.
\]
This formula retains the complete coordinate change and multiplier;
no asymptotic statement about native $E_n$ or $O_n$ follows from
the paper's asymptotic theorem by omitting them.

\section{What the finite estimates do and do not complete}

The parity isometries, Christoffel maps and Gram congruences are
exact representations of the original objects. The nonzero
residual \eqref{eq:residual-lower}, the rational truncation and
error bound \eqref{eq:H-error}, the original observation brackets
\eqref{eq:Q-bracket}--\eqref{eq:relative}, and the explicit
raw-moment comparison \eqref{eq:raw-matrix-bound} are additional
finite calculations and inequalities. Their constants and
receiving maps have all been evaluated above.

They preserve the original complex observed cross entries
\[
 B_i=(\Lambda a_i)^*Q(H)\Lambda v,
 \qquad
 \Phi_{B,N}=\frac{2\Re(\overline{B_1}B_2)}{\omega_{N,k}^{\rm ar}},
\]
with the actual boundary classes, terminal class, and physical
Taylor unit. The bounds certify source/observation approximation
and do not supply an evaluated sign, phase, or required
asymptotic of this last signed product. No such conclusion, and
no imported multiple-orthogonal asymptotic, is claimed.

\begin{thebibliography}{9}
\bibitem{ALM} A. I. Aptekarev, G. L\'opez Lagomasino and
A. Mart\'{\i}nez--Finkelshtein,
\emph{Strong asymptotics for the Pollaczek multiple orthogonal
polynomials ensembles}, arXiv:1410.1261v1 (2014), Introduction:
the two unscaled weights and their positive Stieltjes ratio.
\url{https://arxiv.org/abs/1410.1261v1}.
The original author LaTeX is the source used; its multiple-orthogonal
asymptotic theorem is not assumed for the arithmetic weight.
\bibitem{CAI} Split-Zero programme, \emph{The complete original
action relative to its arithmetic total}, CAI16--18 and its complete original
convolution providers. This is the programme calculation supplying
the actual arithmetic density comparison, not a human-literature
attribution for the Pollaczek weights.
\url{https://github.com/KokunoYumeto/zeta-function-research-reader/blob/2412f4aaa07c3fb19052cc0e537e93c2db9b0b1d/workbenches/splitzero-tandem/continuations/20260920-heat-literature-observed-class/providers/CAI_COMPLETE_PREREQUISITE.tex}.

% BEGIN VERIFIED PUBLIC PROOF LINK
\href{https://github.com/KokunoYumeto/zeta-function-research-reader/blob/487253b5b01851da5b9b98af4e146fecf1fab902/workbenches/splitzero-tandem/continuations/20260919-full-window-source-products/providers/CAI_COMPLETE_PREREQUISITE.tex\#L284}{Source proof CAI16}.
% END VERIFIED PUBLIC PROOF LINK
\bibitem{HG} Split-Zero programme, \emph{Heat growth and original
currents}, HG21 and its displayed proof, retaining every fixed
primary multiplicity and the original Gamma mass.
\url{https://github.com/KokunoYumeto/zeta-function-research-reader/blob/2412f4aaa07c3fb19052cc0e537e93c2db9b0b1d/workbenches/splitzero-tandem/continuations/20260920-heat-literature-observed-class/providers/HEAT_GROWTH_AND_ORIGINAL_CURRENTS.tex}.
The parity recurrences, Green-kernel expansions, finite Schur
inequalities and receiving maps are proved explicitly in this note;
no omitted asymptotic argument is attributed to these programme
providers.

% BEGIN VERIFIED PUBLIC PROOF LINK
\href{https://github.com/KokunoYumeto/zeta-function-research-reader/blob/5992ad50c0f2a06921f1fcaded7b4e38097c0739/workbenches/splitzero-tandem/continuations/20260920-original-kernel-mixed-determinants/providers/HEAT_GROWTH_AND_ORIGINAL_CURRENTS.tex\#L374}{Source proof HG21}.
% END VERIFIED PUBLIC PROOF LINK
\end{thebibliography}

\end{document}
